\section{From safe samples to unobserved rank}\label{sec:transfer}

We continue with the fixed matroid and hidden weight vector, suppressing
them in all conditional expectations.  The next step turns the safe
sample estimate into rank available on labels outside the revealed
mask.  The paths themselves depend on that mask, so the argument needs
a uniform bound on the sets that the paths can leave unspanned.

Write $O=D\cup C$.  For a listed group $h$, let
\[
 \mathcal B_h=\{b\in\mathbb Z:a_h+2\le b\le0,
                              \ b\equiv\varepsilon\pmod 2\},
 \qquad
 P_{h,b}=G_h\cap\bigl(X_h(b)\setminus X_h(b-1)\bigr)
 \quad(b\in\mathcal B_h).
\]
Put $P_{h,b}=\varnothing$ for other endpoints of the chosen parity, and
define the nominal rank statistic
\begin{equation}\label{eq:nominal-rank}
 z_h=\sum_{b\in\mathcal B_h}
 \rk\bigl(P_{h,b}\setminus O\mid X_h(b-2)\cup K_{h,b}\bigr).
\end{equation}
This statistic measures rank before thinning by the focal mask $T_h$
and before replacing $X_h$ by the final path $\widehat F$.  The result
needed for the final accounting is the following estimate.

\begin{lemma}[Sample-to-rank transfer]\label{lem:transfer}
Fix $H$.  For a true group $U_i$, define $z_i=z_h$ and
$\sigma_i=\sigma_h$ if the group is listed with index $h$, and set
both quantities to zero if it is unlisted.  If $n_i\ge\kappa$, then
\begin{equation}\label{eq:rank-transfer}
 \Ex[z_i\mid H]
 \ge\frac{\gamma|Y_i|-2^{-23}n_i}{\kappa},
 \qquad \gamma=\frac\eta{16}=2^{-16}.
\end{equation}
\end{lemma}

We first relate $z_h$ to the safe sampled count $\sigma_h$.  Choose
enough unsacrificed labels to witness the rank in each summand of
\eqref{eq:nominal-rank}.  A safe sample spanned by these witnesses over
its summand's base can be charged to their rank by the density bound.
Every other safe sample belongs to a residual set contained in the
revealed mask.
The main task will be to bound this residual uniformly over the masks
that determine it.

For each summand choose a basis extension
$Z_b\subseteq P_{h,b}\setminus O$ over
$X_h(b-2)\cup K_{h,b}$.  Thus
\begin{equation}\label{eq:rank-witnesses}
 \sum_b|Z_b|=z_h,
 \qquad
 P_{h,b}\setminus O
 \subseteq\cl\bigl(X_h(b-2)\cup K_{h,b}\cup Z_b\bigr).
\end{equation}
The sets $Z_b$ are disjoint, since the birth sets $P_{h,b}$ are
disjoint.  We may choose them by a fixed label order.  Set
\begin{equation}\label{eq:residual-set}
 \mathcal R_h=
 \bigcup_{b\in\mathcal B_h}
 \left(P_{h,b}\setminus
 \cl\bigl(X_h(b-2)\cup K_{h,b}\cup Z_b\bigr)\right).
\end{equation}
By \eqref{eq:rank-witnesses}, $\mathcal R_h\subseteq O$.

There is a deterministic comparison with the safe count from
Section~\ref{sec:safety}:
\begin{equation}\label{eq:safe-residual}
 \sigma_h\le\kappa z_h+|\mathcal R_h|.
\end{equation}
Indeed, fix $b\in\mathcal B_h$ and put
$Q'_b=Q_h(F_{h-1}(b-2)\cup K_{h,b})$.
Since $b-2\ge a_h$, monotonicity of the density map gives
\[
 X_h(b-2)=Q_h(F_{h-1}(b-2))\subseteq Q'_b,
 \qquad K_{h,b}\subseteq Q'_b.
\]
A safe sample at $b$ outside $\mathcal R_h$ belongs to
$D_h\cap(\cl(Q'_b\cup Z_b)\setminus Q'_b)$.
The residual density inequality \eqref{eq:density-residual} bounds the
number of these samples by
$\kappa\rk(Z_b\mid Q'_b)\le\kappa|Z_b|$.
Summing proves \eqref{eq:safe-residual}.

\subsection{Counting marked pivots}

The selected set $\mathcal R_h$ depends on the focal group's mask.
We control it by constructing, before those bits are revealed, a small
family containing every possible such residual with
$z_h\le |G_h|/\kappa$.
The following combinatorial lemma supplies a bound that does not depend
on the number of times at which a path can change.

\begin{lemma}[Marked pivots]\label{lem:pivot-patterns}
Let $M$ be any finite matroid, and let $V$ be a set of $n$ nonloops.
Let $s\ge0$ be an integer.  Fix an ordered sequence of labeled
\emph{old occurrences}, each with a
mark in $\{0,1\}$, whose labels span $E$.  Also fix at most $s$ labeled
\emph{movable occurrences}.  Repetitions of labels are allowed in both
sequences.  Insert the movable occurrences into the old sequence in any
order and positions, and give each movable occurrence either mark.
For $f\in V$, its pivot is the first occurrence after whose insertion
$f$ is spanned; record that occurrence's mark.  The number of resulting
vectors of $n$ marks is at most
\begin{equation}\label{eq:pivot-pattern-bound}
 4^s\sum_{\ell=0}^{s}\binom n\ell,
\end{equation}
where $\binom n\ell=0$ for $\ell>n$.
\end{lemma}

\begin{proof}
We first bound the possible bases retained by greedy processing of an
interleaving.  Remove every old occurrence whose label is spanned by its old
predecessors.  This leaves the span after every old prefix unchanged.
Such an occurrence remains redundant after any insertions, so its
removal preserves all pivot marks.  The surviving old labels form an
ordered basis $b_1,\ldots,b_r$ of $M$.  Write
$B_j=\{b_1,\ldots,b_j\}$, and let $A$ denote all specified movable
occurrences, with rank interpreted through their labels.

For each old position define
\[
 \Delta_j=
 \rk(A\mid B_{j-1})-\rk(A\mid B_j)
 =1-\rk(\{b_j\}\mid B_{j-1}\cup A)\in\{0,1\}.
\]
These numbers are fixed before any insertions are chosen, and
$\sum_j\Delta_j=\rk(A)\le s$.
If $\Delta_j=0$, then $b_j$ increases rank even after all of $A$ has
been adjoined to $B_{j-1}$, and hence it increases rank in every
interleaving.  Thus greedy processing of an interleaving can omit old
occurrences only at the at most $s$ positions with $\Delta_j=1$.
There are at most $2^s$ choices of omitted old occurrences and at most
$2^s$ choices of retained movable occurrences.  Consequently the
retained greedy basis, regarded as a collection of occurrences, has at
most $4^s$ possibilities.

Fix one such basis $J$.  We next identify each pivot from a spanning
subset of $J$, independently of how the other occurrences were inserted.
For every $f\in V$ there is a unique minimal
subset $S_f\subseteq J$ that spans $f$.  To verify uniqueness using only
matroid rank, suppose $S,T\subseteq J$ both span $f$.  Submodularity
applied to $S\cup\{f\}$ and $T\cup\{f\}$ gives
\[
 |S|+|T|\ge |S\cup T|+\rk((S\cap T)\cup\{f\}),
\]
so $S\cap T$ also spans $f$.  Intersecting the finitely many spanning
subsets proves the assertion.  The support $S_f$ is nonempty because
$f$ is a nonloop.  At every prefix, the retained greedy occurrences
span exactly the same flat as the full prefix.  It follows that the
pivot of $f$ is the last occurrence of $S_f$ in the interleaving.

It remains to count the possible marks of these last support members.
For this fixed $J$, assign its retained old occurrences, in their
prescribed order, fixed absolute scores
$2^{s+1},2^{2(s+1)},\ldots$, with positive sign for mark $1$ and
negative sign for mark $0$.  There are at least $s$ unused integer
powers of $2$ before the first old score and between successive old
scores; reserve another $s$ powers after the last.  Every order of the
retained movable occurrences can therefore be represented by assigning
them distinct unused powers in the appropriate gaps, in that order,
with either sign according to their marks.  If no old occurrence is
retained, simply use distinct powers for the movable occurrences.
Every absolute score then exceeds the sum of all preceding absolute
scores.

Let the signed scores of the retained movable occurrences be real
variables, at most $s$ in number.  For each $f$, the sum of scores over
$S_f$ is an affine function of these variables, whose constant term is
the sum of the fixed old scores in $S_f$.  At the representations just
constructed this sum is nonzero, and its sign is precisely the mark of
the last member of $S_f$.  Thus every pivot-mark vector for $J$ is a
strict sign pattern of $n$ affine functions in dimension at most $s$.

The remaining estimate is the classical region bound for affine
hyperplane arrangements; see Schl\"afli~\cite[Section 16, pp.~39--40]{Schlafli1901}.
We include the recurrence and its degenerate cases: $n$ affine functions
in dimension $d$ have at most
$\sum_{\ell=0}^{d}\binom n\ell$ strict sign patterns.  A prescribed
strict pattern defines an open convex region.  Adding a proper affine
hyperplane splits at most as many existing regions as there are regions
of the induced arrangement on that hyperplane, giving the recurrence
$N(n,d)\le N(n-1,d)+N(n-1,d-1)$, with
$N(0,d)=1$ and $N(n,0)\le1$.  A nonzero constant has a fixed sign; an
identically zero function yields no all-strict patterns.  Coincident
hyperplanes and degenerate restrictions cannot increase the recurrence
bound.  Induction gives the displayed binomial sum.  Applying this to
each of the at most $4^s$ bases proves \eqref{eq:pivot-pattern-bound}.
In particular, all movable orders, positions, and marks have already
been counted, with no additional permutation or timeline factor.
\end{proof}

\subsection{Encoding residual sets}

We now apply the counting bound twice to describe the residuals. The
first marked sequence tests whether a label has a valid nominal birth
of the chosen parity. Among labels passing that test, the second tests
whether, in that layer, the label remains
outside the span of the preceding flat, competing test labels, and rank
witnesses. Their intersection is exactly the residual set. Counting both
tests uniformly before exposing the focal group's bits is what permits
a later union bound despite the dependence of the actual residual on
those bits.

\begin{lemma}[A uniform family of residuals]\label{lem:pattern-family}
Fix $H$ and a true group $U_i$ of size $n=n_i\ge\kappa$.
Condition also on the parity and on the $D,C,T$ bits outside $U_i$,
but do not condition on their bits in $U_i$ or on whether $U_i$ is
listed.
There is a deterministic family $\mathfrak R_i\subseteq 2^{U_i}$ with
\begin{equation}\label{eq:pattern-family-size}
 |\mathfrak R_i|\le
 \exp\left(\frac{1000\log\kappa}{\kappa}\,n\right)
\end{equation}
such that, whenever this group is listed and $z_h\le n/\kappa$, every
residual \eqref{eq:residual-set} obtained from witnesses
\eqref{eq:rank-witnesses} belongs to $\mathfrak R_i$.
\end{lemma}

\begin{proof}
If $U_i$ is listed, its index must be one plus the number of listed
lower-weight groups.  Denote this now-fixed index by $h$, whether or
not the focal group is actually listed.  The lower-group paths,
$F(k):=F_{h-1}(k)$, and all $K_{h,b}$ are fixed by the conditioning.
No $D,C,T$ bit on $U_i$ enters this incoming path or these competition
sets.  Write $a=a_h$ and $q=\lfloor n/\kappa\rfloor$.

On any actual listed outcome, Lemma~\ref{lem:generators} supplies a set
$J_h\subseteq U_i$ of at most $q$ density generators and an insertion
time $\tau(g)$ for each $g\in J_h$ such that
\begin{equation}\label{eq:encoded-path}
 X_h(k)=\cl\bigl(F(k)\cup\{g\in J_h:\tau(g)\le k\}\bigr).
\end{equation}
All insertion times can be taken between $a$ and $0$.
When $z_h\le n/\kappa$, the witness union $Z=\bigcup_b Z_b$ has at
most $q$ labels.  We allow every choice of these two subsets of $U_i$,
then use two marked sequences to encode the residual.  An element
chosen for both subsets is treated as two distinct occurrences, one
for each role.

For the first sequence, advance through integer times from $a-2$ to
$1$.  At time $k$ take as old occurrences a fixed label-ordered listing
of $F(k)$; mark all occurrences in this block by $1$ exactly when
$k\in\mathcal B_h$.  These old slots, including their labels, order,
and marks, are fixed by the conditioning and span $E$, since
$F(0)=E$.  For the actual path, insert each $g\in J_h$ at its time
$\tau(g)$, after that time's old block, with the same mark as the
block.  After processing time $k$, the generated flat is exactly $X_h(k)$
by \eqref{eq:encoded-path}.  Hence the set receiving mark $1$ among the
tested labels $U_i$ is precisely
\begin{equation}\label{eq:first-pivot-test}
 \bigcup_{b\in\mathcal B_h} P_{h,b}.
\end{equation}

The second sequence runs through layers of the chosen parity.  Choose
the unique $b_0\in\{a-2,a-1\}$ with
$b_0\equiv\varepsilon\pmod2$, and then take
successive endpoints $b_0+2,b_0+4,\ldots,b_1$, where
$b_1\in\{0,1\}$ has the chosen parity.  The initial span is
$L=X_h(b_0)$.  For each layer $(b-2,b]$ use the following two phases:
\begin{enumerate}[label=\textnormal{(\roman*)},leftmargin=*]
 \item Insert a fixed listing of $K_{h,b}$ as old occurrences with
 mark $0$, followed by the movable witnesses $Z_b$, also with mark $0$.
 Here $Z_b=\varnothing$ if $b\notin\mathcal B_h$.
 \item Insert a fixed listing of $F(b)$ as old occurrences with
 mark $1$, followed by the movable density generators whose insertion
 times satisfy $b-2<\tau(g)\le b$, also with mark $1$.
\end{enumerate}
Again the old slots, their order, and their marks are fixed; their
labels span $E$ by the final endpoint.  Only the density and witness
occurrences move when the focal mask changes.

The two phases are designed to give the following sequence of spans:
\[
 X_h(b-2)
 \xrightarrow[\text{mark }0]{\text{phase (i)}}
 \cl\bigl(X_h(b-2)\cup K_{h,b}\cup Z_b\bigr)
 \xrightarrow[\text{mark }1]{\text{phase (ii)}}
 X_h(b).
\]
We verify these identities before interpreting the pivot marks.
Suppose the generated flat at the preceding endpoint is $X_h(b-2)$.
At the end of phase~\textnormal{(i)} it is exactly
\begin{equation}\label{eq:encoded-interior}
 \cl\bigl(X_h(b-2)\cup K_{h,b}\cup Z_b\bigr).
\end{equation}
The inclusions $K_{h,b}\subseteq F(b)$ and $Z_b\subseteq X_h(b)$,
together with \eqref{eq:encoded-path}, show that the flat at the end
of phase~\textnormal{(ii)} is exactly $X_h(b)$.  This proves the
endpoint assertion by induction from $b_0$.
Density generators originally inserted at $b-1$ have been postponed
to phase~\textnormal{(ii)}.  They affect the intermediate span, but
the endpoints and the explicit span \eqref{eq:encoded-interior}
remain exact.

Now restrict attention to a label passing the first test, with valid
birth $b$.  It is outside $X_h(b-1)$ and hence outside $X_h(b-2)$,
so its second-sequence pivot occurs in this layer.  Its second mark is
$1$ exactly when it is not spanned at the end of phase~\textnormal{(i)}.
By \eqref{eq:encoded-interior}, this is exactly its membership in
$\mathcal R_h$.  Thus the residual is the intersection of the two
sets of labels receiving mark $1$.  Phase~\textnormal{(ii)} can also
span labels born at $b-1$, which explains why the first test is needed:
the second mark alone need not describe a valid birth of the chosen
parity.

We now define $\mathfrak R_i$ without using the focal bits: for every
pair of subsets $J,Z\subseteq U_i$ of size at most $q$, include every
intersection of two marked-pivot sets obtained from the fixed old
sequences just described, allowing arbitrary placements, orders, and
marks of the specified movable occurrences.  The first sequence uses
only the $J$ occurrences; the second uses both roles.  This definition
may include intersections that do not arise from valid paths or
witnesses, which only increases the family.  The preceding construction
shows that it includes every required actual residual.

Let $S_d(n)=\sum_{\ell=0}^{d}\binom n\ell$.
There are at most $S_q(n)^2$ subset choices.  Each sequence has at most
$2q$ movable occurrences, so Lemma~\ref{lem:pivot-patterns} gives
\begin{equation}\label{eq:raw-pattern-count}
 |\mathfrak R_i|
 \le S_q(n)^2\bigl(4^{2q}S_{2q}(n)\bigr)^2.
\end{equation}
This estimate includes all ways of assigning labels to times or
layers.
To bound it with the floors intact, for $0<\alpha\le1/2$ note that
\[
 \alpha^{\alpha n}S_{\lfloor\alpha n\rfloor}(n)
 \le\sum_{\ell=0}^{\lfloor\alpha n\rfloor}
          \binom n\ell\alpha^\ell
 \le(1+\alpha)^n\le e^{\alpha n}.
\]
Use $q\le n/\kappa$ and
$2q\le\lfloor2n/\kappa\rfloor$, with
$\alpha=1/\kappa$ and $\alpha=2/\kappa$, respectively, in
\eqref{eq:raw-pattern-count}.  This yields
\[
 \log|\mathfrak R_i|
 \le\frac n\kappa
       \bigl(6\log\kappa+6+4\log2\bigr)
 \le\frac{1000\log\kappa}{\kappa}\,n,
\]
as claimed.  All labels tested here are nonloops by the initial
candidate filter, and the argument uses no representability property
of the matroid.
\end{proof}

\subsection{Completing the sample-to-rank estimate}

The residual family is fixed before the focal mask is exposed.  We can
therefore bound the chance that any large member lies entirely in $O$,
even though the actual residual is chosen using that mask.

\begin{proof}[Proof of Lemma~\ref{lem:transfer}]
Put $n=n_i$.  On the unlisted branch $D\cap U_i$ is empty, so
Lemma~\ref{lem:safe}, with both sides extended by zero on that branch,
can be averaged over $D$.  Since each label belongs to $D$ with
probability $1/4$, it gives
\begin{equation}\label{eq:averaged-safe-count}
 \Ex[\sigma_i\mid H]
 \ge\frac\eta{16}|Y_i|-\frac{n}{4\kappa}
 \ge\gamma|Y_i|-\frac n\kappa.
\end{equation}

Use the conditioning in Lemma~\ref{lem:pattern-family}.  Within $U_i$,
the events $\{e\in O\}$ remain independent across labels, each with
probability
\[
 p_O=1-(1-1/4)(1-t)=\frac14+\frac34t\le\frac12.
\]
Thus, for each fixed $R_0\in\mathfrak R_i$,
$\Prob(R_0\subseteq O)\le2^{-|R_0|}$ under this conditioning.
Let $\delta=2^{-26}$ and let $\mathcal E_i$ be the event that some
$R_0\in\mathfrak R_i$ satisfies $R_0\subseteq O$ and $|R_0|>\delta n$.
The union bound, without conditioning on listing, gives
\begin{align}
 \Prob(\mathcal E_i\mid H,\varepsilon,
                  \text{masks outside }U_i)
 &\le\exp\left(\frac{1000\log\kappa}{\kappa}n
                         -\delta n\log2\right)\notag\\
 &=2^{-(2^{-26}-100000\,2^{-100})n}
 \le2^{-2^{-27}n}
 \le2^{-2^{73}}<2^{-26}.
 \label{eq:residual-tail}
\end{align}
Here we used $\kappa=2^{100}$ and $n\ge\kappa$; in particular the
bound already holds at the smallest analyzed group size.

Define an error variable $e_i$ as follows.  On a listed outcome with
$z_i\le n/\kappa$, choose the witnesses in
\eqref{eq:rank-witnesses} and set $e_i=|\mathcal R_h|$.
Set $e_i=0$ on all other outcomes.  In its nonzero branch,
$\mathcal R_h\in\mathfrak R_i$ and $\mathcal R_h\subseteq O$, so
\[
 e_i\le\delta n+n\ind{\mathcal E_i},
 \qquad
 \Ex[e_i\mid H]\le2\delta n=2^{-25}n.
\]
The latter bound follows first under the conditioning of
\eqref{eq:residual-tail}, then by averaging that conditioning away.
No independence is asserted for the selected residual itself.

In every outcome,
\[
 \sigma_i\le\kappa z_i+e_i.
\]
For listed outcomes with $z_i\le n/\kappa$ this is
\eqref{eq:safe-residual}; for listed outcomes with $z_i>n/\kappa$
it follows from $\sigma_i\le n<\kappa z_i$; for unlisted outcomes
both statistics vanish.  Taking expectations and using
\eqref{eq:averaged-safe-count} gives
\[
 \kappa\Ex[z_i\mid H]
 \ge\gamma|Y_i|-(\kappa^{-1}+2^{-25})n
 \ge\gamma|Y_i|-2^{-23}n,
\]
which proves \eqref{eq:rank-transfer}.
\end{proof}
